\section{El ecosistema de Agent: de AgentRank a OS Agent}

La sección anterior trató el contenedor de producto de YouWare; esta sección pasa a la valoración del ecosistema de Agent. Ming Chaoping (明超平) compara el ecosistema de Agent con la sociedad humana: puede haber un OS Agent que coordine a todos los Agent verticales, o bien un conjunto de Agent verticales que se invocan entre sí. En el futuro no importará solo si un Agent individual puede completar una tarea, sino quién puede ser invocado, en quién se confía y quién logra entrar en el efecto de red.

\begin{figure}[H]
\centering
\includegraphics[width=0.94\textwidth]{agent-ecosystem-two-models.png}
\caption{Dos ecosistemas de Agent: pequeño y fuerte al estilo de Singapur frente a grande y completo al estilo de Estados Unidos. Diagrama conceptual propio, elaborado a partir de la conversación entre 01:37:12 y 01:40:44.}
\end{figure}

\begin{knowledgebox}{Lectura de la figura: el ecosistema tiene dos formas de organización}
A la izquierda, Singapore-like representa lo pequeño y fuerte, con conexiones abiertas y dependencia del ecosistema externo; a la derecha, US-like representa lo grande y completo, autosuficiente y con capacidades de plataforma completas. Una empresa de aplicaciones debe decidir primero si se convertirá en un nodo de capacidad vertical o si buscará ser el coordinador a nivel de OS.
\end{knowledgebox}

\subsection{De PageRank a AgentRank}

La subsección anterior trató las formas del ecosistema de Agent; esta examina la clasificación y la distribución. En la era de la web, PageRank mide la importancia de una página mediante sus relaciones de enlace. En la era de los Agent, un mecanismo similar podría convertirse en AgentRank: las relaciones de invocación entre usuarios, OS Agent, Agent verticales, herramientas y contenidos formarán nuevos mecanismos de clasificación y distribución. Quien sea invocado por más nodos de alta confianza podrá obtener más tareas, datos e ingresos.

\begin{figure}[H]
\centering
\includegraphics[width=0.96\textwidth]{page-rank-to-agent-rank.png}
\caption{De PageRank a AgentRank: las páginas se clasifican por enlaces; los Agent podrían clasificarse por invocaciones y confianza. Diagrama conceptual propio, elaborado a partir de la conversación entre 01:40:44 y 01:44:02.}
\end{figure}

\begin{knowledgebox}{Lectura de la figura: los enlaces se vuelven invocaciones y el peso se vuelve confianza}
En la figura, Page, Links y PageRank se trasladan a Agent y AgentRank. La «clasificación» en el ecosistema de Agent no se limita a los resultados de búsqueda: combina la distribución de tareas, la invocación de herramientas, la autenticación de identidad y la confiabilidad de los resultados.
\end{knowledgebox}

\begin{knowledgebox}{Términos clave: la red de Agent}
\noindent\begin{tabularx}{\textwidth}{p{0.18\textwidth}p{0.36\textwidth}X}
\toprule
Término & Significado & Función en este episodio \\
\midrule
OS Agent & Agente de nivel superior que coordina a otros Agent y al entorno & Podría convertirse en el punto de entrada de la próxima generación. \\
Vertical Agent & Agente dedicado a una tarea vertical específica & Podría convertirse en un nodo de capacidad coordinado por otros. \\
AgentRank & Clasificación basada en invocaciones, confianza y resultados & Determina la distribución del tráfico y de las tareas. \\
Network Effect & Cuantos más usuarios y Agent entran, mayor es el valor & Una de las barreras de largo plazo de una empresa de aplicaciones. \\
Identity & Identidad, permisos y prueba de origen de un Agent & Resuelve problemas de confianza y seguridad. \\
\bottomrule
\end{tabularx}
\end{knowledgebox}

\subsection{Efecto de red y problema de confianza entre Agent}

Cuando los Agent se invocan entre sí, aparece el efecto de red, pero también aparece el problema de la confianza. Ming Chaoping menciona que en el futuro los Agent formarán tribus, como en la sociedad humana, y necesitarán documentos de identidad, cerraduras con contraseña y gestión de la confianza. Un Agent no puede invocar a otro Agent en nombre del usuario de forma arbitraria, ni entregar sin control los datos del usuario o ejecutar operaciones de alto riesgo. La identidad, los permisos, la auditoría y la confianza se convertirán en la infraestructura del ecosistema de Agent.

\begin{figure}[H]
\centering
\includegraphics[width=0.92\textwidth]{agent-network-effect.png}
\caption{Efecto de red de los Agent: las invocaciones, la identidad, la confianza y las herramientas entre Agent forman un ecosistema. Diagrama conceptual propio, elaborado a partir de la conversación entre 01:44:02 y 02:03:33.}
\end{figure}

\begin{knowledgebox}{Lectura de la figura: el efecto de red y las restricciones de seguridad aparecen al mismo tiempo}
En el centro de la figura está Agent Network, rodeado por Identity, Trust, Tools, Users, Calls y Rank. Cuanto más valiosa es la red, más necesita identidad y permisos; de lo contrario, las relaciones de invocación se convierten en superficie de ataque.
\end{knowledgebox}

\noindent\begin{tabularx}{\textwidth}{p{0.18\textwidth}p{0.36\textwidth}X}
\toprule
Problema de gobernanza & Por qué aparece & Posibles mecanismos de producto \\
\midrule
Identidad & El Agent debe actuar en nombre de un usuario o una organización & Cuentas, firmas, permisos, credenciales. \\
Confianza & Los Agent se invocan y se recomiendan entre sí & Historial de invocaciones, calificaciones, auditoría, prueba de origen. \\
Permisos & No todos los Agent pueden acceder a todos los datos & Mínimo privilegio, confirmación del usuario, aislamiento de alcance. \\
Responsabilidad & Tras un error hay que saber quién lo desencadenó y quién lo ejecutó & Registros, reproducción, flujos de trabajo rastreables. \\
Clasificación & Hay demasiados nodos de capacidad y se necesita un mecanismo de distribución & AgentRank, asignación de tareas, recomendación según el contexto. \\
\bottomrule
\end{tabularx}

\subsection{Los Agent de hoy son como quien acaba de tomar un palo para el fuego}

Esta subsección explica la metáfora del título. Ming Chaoping dice que los Agent de hoy se parecen a un gorila que acaba de tomar una piedra o un palo para el fuego: muy toscos, propensos a romper cosas, pero con una dirección de enorme alcance. La metáfora subraya dos hechos: por un lado, los Agent de hoy aún no son estables, y el uso de herramientas, el contexto, el entorno, la identidad y la confianza son todavía muy primitivos; por otro, en cuanto las herramientas y el entorno maduren, la forma de producir se reorganizará con rapidez.

\begin{figure}[H]
\centering
\includegraphics[width=0.94\textwidth]{ape-with-fire-stick.png}
\caption{La etapa actual de los Agent: como un gorila que acaba de tomar un palo para el fuego, con capacidades toscas pero una dirección de enorme alcance. Diagrama conceptual propio, elaborado a partir de la conversación entre 01:46:30 y 01:49:50.}
\end{figure}

\begin{knowledgebox}{Lectura de la figura: Primitive y Potential deben verse al mismo tiempo}
A la izquierda, Primitive significa tosco, con muchos intentos fallidos y propenso a romper cosas; a la derecha, Potential significa que la forma de producir se reorganiza cuando las herramientas mejoran. Al leer la figura no hay que limitarse a burlarse de la inestabilidad actual, pero tampoco hay que pasar por alto los límites actuales de las capacidades.
\end{knowledgebox}

\subsection{Resumen de la sección}

La oportunidad de largo plazo de las aplicaciones de Agent no está en herramientas aisladas, sino en el ecosistema, la clasificación, la identidad y el entorno. Lo que YouWare quiere construir no es solo una herramienta de coding, sino un entorno que pueda sostener la creación y la relación con el OS Agent.
